\tocsection{\textbf{11/3/26 Lecture 20:} Likelihood Ratio Tests and the Neyman-Pearson Lemma}

\begin{example}{\textbf{Example:} $z$-test}
    Suppose we have $X$ and a summary statistic $Z(X) \sim N(\theta, 1)$. Define the upper $\alpha$ quantile
    \[ z_{\alpha} = \Phi^{-1}(1 - \alpha) \quad \tt{where } \Phi \tt{ is the } N(0, 1) \tt{ c.d.f.} \]
    \underline{1-sided $z$-test}: $H_0 \col \theta \leq 0$ vs.\ $H_1 \col \theta > 0$ and our test function is
    \[ \phi_1(X) = 1\{Z(X) > z_{\alpha}\} \]
    \underline{2-sided $z$-test}: $H_0 \col \theta = 0$ vs.\ $H_1 \col \theta \neq 0$ with test function
    \[ \phi_2(X) = 1\{\abs{Z(X)} > z_{\alpha/2}\} \]
    Note that $\phi_1$ rejects for large values of $Z(X)$ and $\phi_2$ rejects for large values of $\abs{Z(X)}$.
    \fig{0.4}{Lecture20Fig1}{}
    \textbf{Power:}
    \fig{0.45}{Lecture20Fig2}{}
    If we are testing $H_0 \col \theta \leq 0$, $\beta_{\phi_1}$ is valid but not $\beta_{\phi_2}$.\\

    If we are testing $H_0 \col \theta = 0$, both $\beta_{\phi_1}$ and $\beta_{\phi_2}$ are valid, but $\beta_{\phi_2}$ is better for $\theta < 0$.
\end{example}

It turns out $\phi_1$ is the best test (in the sense it uniformly maximizes $\beta_{\phi}$ on $H_1$ while still being level-$\alpha$). How do we show this?

\subsection*{Likelihood Ratio Test}
\sectionspace

Start with simple vs.\ simple:
\[ H_0 \col X \sim P_0 \qquad H_1 \col X \sim P_1 \]
Assume WLOG $P_0$ and $P_1$ have densities $p_0, p_1$ with respect to a dominating measure $\mu$ (they are absolutely continuous with respect to $P_0 + P_1$).

\begin{definition}[\textbf{Definition 20.1:} Likelihood Ratio]
    The \emph{likelihood ratio} is
    \[ \tt{LR}(X) = \frac{p_1(X)}{p_0(X)} \]
    with $\tt{LR}(X) = \infty$ if $p_1 > 0$ and $p_0 = 0$ (and define it to be whatever if they are both $0$).
\end{definition}

\begin{definition}[\textbf{Definition 20.2:} Likelihood Ratio Test]
    The \emph{likelihood ratio test} is
    \[ \phi_{\alpha}^{*}(X) = \left\{\begin{matrix} 1 & \tt{LR}(X) > c_{\alpha} \\ \gamma_{\alpha} & \tt{LR}(X) = c_{\alpha} \\ 0 & \tt{LR}(X) < c_{\alpha} \end{matrix}\right. \]
    where $c_{\alpha}, \gamma_{\alpha}$ are chosen to make $\E_0\phi_{\alpha}^{*}(X) = \alpha$.
\end{definition}

\textbf{Intuition for LRT:}
\begin{align*}
    \tt{Power under } H_1 &\col \max_{R} \int_{R} p_1(x) d\mu(x)\\
    \tt{Significance budget} &\col \int_{R} p_0(x) d\mu(x) \leq \alpha
\end{align*}

\begin{theorem}[\textbf{Theorem 20.3:} Neyman-Pearson Lemma]
    $\phi_{\alpha}^{*}$ is optimal among level-$alpha$ tests for testing $H_0 \col X \sim P_0$ vs.\ $H_1 \col X \sim P_1$.
\end{theorem}
\begin{proof}{Proof:}
    We want to show $\phi_{\alpha}^{*}$ solves
    \[ \underset{\phi \col \cal{X} \to [0, 1]}{\tt{maximize}}\ \E_1[\phi(X)] \quad \tt{subject to} \quad \E_0[\phi(X)] \leq \alpha \]
    In Lagrange form,
    \begin{align*}
        \underset{\phi \col \cal{X} \to [0, 1]}{\tt{maximize}}\left(\E_1[\phi(X)] - \lambda\E_0\phi(X)\right) &= \int_{\cal{X}} \phi(x)(p_1(x) - \lambda p_0(x)) d\mu\\
        &= \int_{\{p_0 > 0\}} \phi(x)\left(\frac{p_1(x)}{p_0(x)} - \lambda\right) dP_0(x) + \int_{\{p_0 = 0\}} \phi(x)p_1(x) d\mu
    \end{align*}.
    To maximize this, we want $\phi(x) = 1$ on $\{p_0 = 0\}$, and $\phi(x) = 1$ if $\tt{LR}(x) > \lambda$, and $\phi(x) = 0$ if $\tt{LR}(x) < \lambda$, so
    \[ \phi(x) = \left\{\begin{matrix} 1 & \tt{LR} > \lambda \\ 0 & \tt{LR} < \lambda \\ \tt{arbitrary} & \tt{LR} = \lambda \end{matrix}\right. \]
    Thus, $\phi_{\alpha}^{*}$ maximizes the Lagrangian if we take $\lambda = c_{\alpha}$.\\

    For any other $\tilde{\phi}$ with $\E_0\tilde{\phi}(X) \leq \alpha$,
    \begin{align*}
        \E_1[\tilde{\phi}(X)] &\leq \E_1[\tilde{\phi}] - c_{\alpha}\E_0[\tilde{\phi}] + c_{\alpha}\alpha\\
        &\leq \E_1[\phi_{\alpha}^{*}] - c_{\alpha}\E_0[\phi_{\alpha}^{*}] + c_{\alpha}\alpha\\
        &= \E_1\phi_{\alpha}^{*}
    \end{align*}

    $\qed$
\end{proof}

\begin{example}{\textbf{Example:} Binomial}
    $X \sim \Binom(n, \theta)$,
    \[ H_0 \col \theta = 0.5 \quad \tt{vs} \quad H_1 \col \theta = 0.51 \quad \tt{at level } \alpha = 0.05 \]
    We have
    \[ p_{\theta}(x) = \binom{n}{x}\theta^{x}(1 - \theta)^{n - x} \]
    so
    \[ \tt{LR}(X) = \frac{0.51^{X}0.49^{n - X}}{0.5^{n}} \]
    Since we reject for large LR, we just need to consider $0.51^{x}0.49^{-x} = \left(\frac{0.51}{0.49}\right)^{x}$, and as this is increasing in $x$, we just need to reject when $x$ is large.\\

    Thus, our test statistic is just $X$ and we need to choose a threshold $c_0$. Thus, $c_0$ should be the lowest 95th percentile of $\Binom(n, 0.5)$. That is,
    \[ c_0 = \min\{c \col \bb{P}_{\theta_0}(X > c) \leq \alpha\} \]
    Since $\bb{P}_{\theta_0}(X > c_0) < \alpha$ as $0.05$ is not dyadic, we can set
    \[ \gamma_{\alpha} = \frac{\alpha - \bb{P}_{\theta_0}(X > c_0)}{\bb{P}_{\theta_0}(X = c_0)} \]
    and we get the test
    \[ \phi(X) = \left\{\begin{matrix} 1 & X > c_0 \\ \gamma_{\alpha} & X = c_0 \\ 0 & X < c_0 \end{matrix}\right. \]
    Note that as all of our choices were independent of $H_1$, it also works for $H_0 \col \theta = 0.5$, $H_1 \col \theta = 0.52$, and generally, this would be the test for
    \[ H_0 \col \theta = 0.5 \qquad H_1 \col \theta > 0.5 \]
\end{example}
